% GENERATED by tools/preliminary_extraction_table.py --- do not hand-edit.
% Numbers: data/gap/archive/v0_4_extraction_run.json (computed from the v0.4 GAP
%   on archive/attacker-profiling and from data/gap/gap_v0.5.json).
% Wording: data/gap/archive/preliminary_extraction_labels.json.
% Caption SESSION-DRAFTED --- flagged for the voice pass.
% Requires booktabs (already in the preamble).

\begin{table}[htbp]
\centering
\footnotesize
\setlength{\tabcolsep}{4pt}
\caption[Modes of obtaining tactic dependency that were tried]{The modes of obtaining dependency between tactics that were attempted, and what each returned. \emph{Direction from} is the column the entry turns on: a mode that takes its direction from an assumed tactic ordering cannot learn an ordering, and can express no loop at all, since a total order admits no backward edge --- where the graph that ships preserves loops deliberately. \emph{Tactic transitions} counts distinct ordered transitions between two different tactics, of the 182 that 14 tactics admit, with techniques collapsed onto their primary tactic --- the granularity the model works at. All three rows are measured on a single build, the superseded v0.4 profile of 2026-04-19 (ATT\&CK Enterprise, 216 parent techniques): it is the only artefact in which the two abandoned modes were ever merged, and the shared build is what makes the final column a comparison between modes rather than between corpus vintages. Carried forward alone, Attack Flow reaches 113 of the 182 transitions (62\,\%) in the graph this dissertation uses --- 478 technique edges over 38 incidents (Attack Flow corpus v3.1.1, ATT\&CK Enterprise v19.1) --- quoted over the same 182 transitions, with the split that model makes within defence evasion folded back so the two figures share an axis. No threshold is reported anywhere in the table: the mining parameters were varied throughout the exploratory period, and the table is built on what does not move with them. Not attempted, and so absent here: Natural-language extraction from threat-intelligence prose; Process mining over incident event logs --- ruled out on the literature and never run, named here so the table cannot be read as reporting them.}
\label{tab:preliminary-extraction}
\begin{tabular}{@{}l >{\raggedright\arraybackslash}p{0.18\textwidth} >{\raggedright\arraybackslash}p{0.12\textwidth} r >{\raggedright\arraybackslash}p{0.29\textwidth}@{}}
\toprule
Mode & A dependency rests on & Direction from & \shortstack[r]{Tactic\\transitions} & Why it was not used \\
\midrule
Co-occurrence mining & two techniques being used by the same actor or tool & assumed tactic ordering & 19 (10\,\%) & Association is symmetric, so the ordering had to be assumed rather than learned; and because a total order admits no backward edge, no loop could be expressed. \\
Keyword extraction & precondition wording in a technique's own description & assumed tactic ordering & 25 (14\,\%) & A description that names another technique is cross-referencing it, not asserting a dependency on it; direction again came from the assumption. \\
\textbf{Attack Flow} & a dependency an analyst drew in one incident & the analyst's drawing & 91 (50\,\%) & Adopted. \\
\bottomrule
\end{tabular}
\end{table}
